\documentclass[aspectratio=169,11pt]{beamer}
% Shared Beamer style for practice contest solution walkthroughs.
\usepackage[utf8]{inputenc}
\usepackage[T1]{fontenc}
\usepackage{lmodern}
\usepackage{amsmath,amssymb}
\usepackage{xcolor}
\usepackage{hyperref}
\usepackage{listings}

\definecolor{CPNavy}{HTML}{172033}
\definecolor{CPBlue}{HTML}{2563EB}
\definecolor{CPGreen}{HTML}{16A34A}
\definecolor{CPOrange}{HTML}{F97316}
\definecolor{CPGray}{HTML}{64748B}
\definecolor{CPLight}{HTML}{F8FAFC}
\definecolor{CPCodeBg}{HTML}{F1F5F9}
\definecolor{CPCodeRule}{HTML}{CBD5E1}

\usetheme{default}
\usefonttheme{professionalfonts}
\setbeamertemplate{navigation symbols}{}
\setbeamersize{text margin left=0.65cm,text margin right=0.65cm}
\setbeamertemplate{itemize item}{\color{CPBlue}$\blacktriangleright$}
\setbeamertemplate{itemize subitem}{\color{CPGray}$\triangleright$}

\setbeamercolor{normal text}{fg=CPNavy,bg=white}
\setbeamercolor{frametitle}{fg=white,bg=CPNavy}
\setbeamercolor{title}{fg=white,bg=CPNavy}
\setbeamercolor{subtitle}{fg=CPLight,bg=CPNavy}
\hypersetup{colorlinks=true,urlcolor=CPBlue}

\setbeamertemplate{frametitle}{%
  \nointerlineskip%
  \begin{beamercolorbox}[wd=\paperwidth,ht=0.88cm,dp=0.22cm,leftskip=0.55cm,rightskip=0.35cm]{frametitle}
    \usebeamerfont{frametitle}\insertframetitle\hfill\small\insertframenumber/\inserttotalframenumber
  \end{beamercolorbox}%
}

\setbeamertemplate{footline}{%
  \leavevmode%
  \hbox{%
    \begin{beamercolorbox}[wd=.58\paperwidth,ht=2.4ex,dp=1ex,leftskip=.5cm]{author in head/foot}%
      \color{CPGray}\insertshorttitle
    \end{beamercolorbox}%
    \begin{beamercolorbox}[wd=.42\paperwidth,ht=2.4ex,dp=1ex,rightskip=.5cm plus1fil]{date in head/foot}%
      \hfill\color{CPGray}\insertshortauthor
    \end{beamercolorbox}%
  }%
}

\setbeamertemplate{title page}{%
  \vspace*{1.25cm}
  \begin{beamercolorbox}[wd=\paperwidth,sep=0.65cm,leftskip=0.15cm,rightskip=0.15cm]{title}
    {\color{white}\Huge\bfseries\inserttitle\par}
    \vspace{0.35cm}
    {\color{CPLight}\Large\insertsubtitle\par}
    \vspace{0.6cm}
    {\color{CPLight}\large\insertauthor\par}
  \end{beamercolorbox}
  \vfill
}

\newcommand{\sectiontag}[1]{\textbf{\color{CPBlue}#1}}
\newenvironment{tightitemize}{\begin{itemize}\small\setlength{\itemsep}{0.18cm}}{\end{itemize}}
\lstdefinestyle{cpcode}{
  language=Python,
  basicstyle=\ttfamily\tiny,
  keywordstyle=\color{CPBlue}\bfseries,
  commentstyle=\color{CPGray}\itshape,
  stringstyle=\color{CPGreen},
  backgroundcolor=\color{CPCodeBg},
  frame=single,
  framerule=0.25pt,
  rulecolor=\color{CPCodeRule},
  showstringspaces=false,
  columns=fullflexible,
  keepspaces=true,
  breaklines=true,
  breakatwhitespace=false,
  tabsize=4,
  xleftmargin=0.08cm,
  xrightmargin=0.08cm,
  aboveskip=0.08cm,
  belowskip=0.08cm
}


\title[grid]{Grid}
\subtitle{Day 4 solution walkthrough}
\author[Solution Deck]{Competitive Programming Bootcamp}
\date{}

\begin{document}

\begin{frame}[plain]
  \titlepage
\end{frame}

% BEGIN GENERATED STATEMENT INTERPRETATION
\begin{frame}{Interpret the Word Problem}
\small
\sectiontag{Statement excerpts:}
\begin{tightitemize}
  \item \textit{``jumping exactly \$k\$ squares''}
  \item \textit{``minimum number of moves required''}
\end{tightitemize}

\vspace{0.18cm}
\sectiontag{Programming interpretation:}
\begin{tightitemize}
  \item Each cell is a node and each legal jump is an unweighted edge.
  \item Because every jump costs one move, BFS gives the shortest number of moves.
  \item Unvisited cells keep distance -1, which also handles unreachable targets.
\end{tightitemize}
\end{frame}
% END GENERATED STATEMENT INTERPRETATION







\begin{frame}{Problem Model}
\small
\sectiontag{Textbook connection:} CP4 4.4 a. On Unweighted Graph: BFS

\vspace{0.25cm}
\sectiontag{Core technique:} BFS on a jump grid

\vspace{0.25cm}
\begin{tightitemize}
  \item Each cell is a node.
  \item The digit in a cell tells the exact jump length.
  \item Each legal jump costs one move.
\end{tightitemize}
\end{frame}

\begin{frame}{How to Recognize the Approach}
\begin{tightitemize}
  \item All moves have equal cost, so BFS finds the minimum number of jumps.
  \item The graph is implicit; neighbors are computed from the current digit.
  \item A distance grid records both visited status and shortest distance.
\end{tightitemize}
\end{frame}

\begin{frame}{Algorithm}
\begin{tightitemize}
  \item Parse the grid as integers.
  \item Set distance[0][0] = 0 and push the start cell into a queue.
  \item Pop cells in BFS order.
  \item For each of four directions, jump exactly d cells.
  \item If the target cell is in bounds and unvisited, assign distance + 1 and enqueue it.
\end{tightitemize}
\end{frame}

\begin{frame}{Walkthrough of the Key State}
\begin{tightitemize}
  \item A cell with value 3 has at most four outgoing moves: up/down/left/right by 3.
  \item Visited cells are not enqueued again because BFS already found their shortest path.
  \item The answer is distance at the bottom-right cell, or -1 if unreached.
\end{tightitemize}
\end{frame}

\begin{frame}{Why the Algorithm Is Correct}
\begin{tightitemize}
  \item BFS processes states in nondecreasing number of moves.
  \item The first time a cell is reached, no shorter path to it can exist.
  \item The generated moves are exactly the legal jumps from each cell.
\end{tightitemize}
\end{frame}

\begin{frame}{Complexity and Implementation Notes}
\small
\sectiontag{Complexity:} O(RC) time and O(RC) memory.

\vspace{0.3cm}
\begin{tightitemize}
  \item Decode each input row before converting characters to integers.
  \item Jump length 0 creates no useful new move unless already at the destination.
\end{tightitemize}
\end{frame}



% BEGIN GENERATED CODE WALKTHROUGH
\begin{frame}[fragile]{Code: Read Digits as Integers}
\scriptsize
\sectiontag{Source lines:} \texttt{grid.py:42--60}

\vspace{0.08cm}
\begin{columns}[T,totalwidth=\textwidth]
\begin{column}{0.58\textwidth}
\begin{lstlisting}[style=cpcode]
import sys
from collections import deque

def main():
    input = sys.stdin.buffer.readline

    rows, cols = map(int, input().split())

    grid = []

    for _ in range(rows):
        row = input().strip().decode()

        grid.append([int(char) for char in row])
\end{lstlisting}
\end{column}
\begin{column}{0.38\textwidth}
\sectiontag{What this chunk does}
\begin{tightitemize}
  \item The grid rows arrive as compact digit strings.
  \item Each character is converted into the jump length for that cell.
  \item The resulting 2D list supports constant-time access during BFS.
\end{tightitemize}
\end{column}
\end{columns}
\end{frame}

\begin{frame}[fragile]{Code: Initialize BFS}
\scriptsize
\sectiontag{Source lines:} \texttt{grid.py:62--73}

\vspace{0.08cm}
\begin{columns}[T,totalwidth=\textwidth]
\begin{column}{0.58\textwidth}
\begin{lstlisting}[style=cpcode]
distance = [[-1] * cols for _ in range(rows)]

queue = deque([(0, 0)])
distance[0][0] = 0

directions = [(-1, 0), (1, 0), (0, -1), (0, 1)]  # up  # down  # left  # right
\end{lstlisting}
\end{column}
\begin{column}{0.38\textwidth}
\sectiontag{What this chunk does}
\begin{tightitemize}
  \item distance stores the best known move count for each cell.
  \item The start cell is the top-left corner with distance zero.
  \item The directions list covers the four cardinal jumps.
\end{tightitemize}
\end{column}
\end{columns}
\end{frame}

\begin{frame}[fragile]{Code: Explore Jump Edges}
\scriptsize
\sectiontag{Source lines:} \texttt{grid.py:75--99}

\vspace{0.08cm}
\begin{columns}[T,totalwidth=\textwidth]
\begin{column}{0.58\textwidth}
\begin{lstlisting}[style=cpcode]
while queue:
    row, col = queue.popleft()

    jump = grid[row][col]

    for dr, dc in directions:
        new_row = row + dr * jump
        new_col = col + dc * jump

        if new_row < 0 or new_row >= rows or new_col < 0 or new_col >= cols:
            continue

        if distance[new_row][new_col] != -1:
            continue

        distance[new_row][new_col] = distance[row][col] + 1

        queue.append((new_row, new_col))
\end{lstlisting}
\end{column}
\begin{column}{0.38\textwidth}
\sectiontag{What this chunk does}
\begin{tightitemize}
  \item BFS removes cells in increasing move count.
  \item The cell's digit determines the exact jump length.
  \item A cell is queued only the first time it is reached, which is its shortest distance.
\end{tightitemize}
\end{column}
\end{columns}
\end{frame}

\begin{frame}[fragile]{Code: Return Target Distance}
\scriptsize
\sectiontag{Source lines:} \texttt{grid.py:101--106}

\vspace{0.08cm}
\begin{columns}[T,totalwidth=\textwidth]
\begin{column}{0.58\textwidth}
\begin{lstlisting}[style=cpcode]
    print(distance[rows - 1][cols - 1])

main()
\end{lstlisting}
\end{column}
\begin{column}{0.38\textwidth}
\sectiontag{What this chunk does}
\begin{tightitemize}
  \item The bottom-right distance is the required answer.
  \item If BFS never reached it, the value remains -1.
\end{tightitemize}
\end{column}
\end{columns}
\end{frame}
% END GENERATED CODE WALKTHROUGH

\begin{frame}{Source}
\small
\sectiontag{Solution file:} \texttt{practice\_contests/day\_4/grid.py}

\vspace{0.3cm}
\sectiontag{Problem:} \url{https://open.kattis.com/problems/grid}

\vspace{0.45cm}
Use this deck with the source code open beside it. The slides explain the reasoning; the code shows the exact input handling and output formatting.
\end{frame}

\end{document}
